\chapter{A History of Stability-Related Incidents}
\label{ch:history}

Chernobyl was not the first time reactor dynamics asserted themselves. A series of
earlier events, some deliberate experiments and some accidents, mapped out the
behaviour of reactors near the edge of stability and shaped the safety principles
that the RBMK, tragically, did not fully embody. This chapter surveys the most
instructive cases as context for the Chernobyl analysis to follow.

\section{Deliberate destructive tests: BORAX and SPERT}

In the 1950s and 1960s the United States conducted a remarkable series of
\emph{intentional} excursion experiments to learn how reactors behave when pushed
past prompt criticality. The BORAX and SPERT (Special Power Excursion Reactor
Test) programmes inserted large step reactivities into water-moderated reactors
and measured the response. The central finding validated the Nordheim--Fuchs
picture: in an under-moderated water reactor, the prompt negative void and Doppler
feedback produced a self-limiting power spike. As the water heated and boiled, the
negative void coefficient sharply reduced reactivity and terminated the
excursion---often before any mechanical action. The BORAX-I reactor was eventually
destroyed in a deliberately oversized excursion, confirming that even when the
energy release fragments the core, the feedback that ends the transient is the
machine's own physics. These tests are the empirical foundation for the claim that
a negative-void, negative-Doppler reactor is inherently self-limiting.

\section{SL-1 (1961): a prompt-critical fatal excursion}

The SL-1 (Stationary Low-Power Reactor No.~1) accident in Idaho on 3~January 1961
was the first and only fatal reactor accident in U.S.\ history. A single central
control rod was manually withdrawn far beyond its normal range---by hand, during
maintenance---inserting a large prompt-critical reactivity in an instant. The
power rose to perhaps $20\unit{GW}$ within milliseconds; the water flashed to steam
and the resulting pressure surge (a water-hammer) lifted the entire reactor vessel
several feet, killing the three operators. SL-1 demonstrated, lethally, two
themes that recur at Chernobyl: a single excessive reactivity insertion can take a
reactor prompt-critical faster than anyone can react, and the destructive energy
is delivered not directly by neutrons but by the resulting steam explosion.

\section{NRX (1952) and the Windscale fire (1957)}

Canada's NRX reactor suffered a partial meltdown in 1952 after a combination of
operator error and control-rod malfunction led to a power excursion with
inadequate shutdown; it became an early lesson in the importance of reliable,
fast, and unambiguous shutdown systems. The 1957 Windscale fire in the United
Kingdom, though primarily a graphite-fire and Wigner-energy event rather than a
prompt excursion, underscored the hazards of graphite-moderated, gas-cooled
reactors and of stored energy in graphite---issues that resonate with the RBMK's
graphite moderator.

\section{Fermi-1 (1966): partial meltdown in a fast reactor}

The Enrico Fermi-1 sodium-cooled fast breeder reactor suffered a partial fuel
meltdown in 1966 when a loose flow-blocking plate starved several fuel assemblies
of coolant. Although triggered by a flow blockage rather than a reactivity event,
Fermi-1 highlighted the tight coupling between thermal-hydraulics and core
integrity, and the consequences of local cooling failure---an enduring theme in
thermal safety.

\section{Boiling-water reactor instability events}

Boiling-water reactors carry a designed-in stability concern because their void
feedback is transported by the coolant and therefore phase-lagged
(Chapter~\ref{ch:linear}). Several plants experienced unintended power
oscillations in low-flow, high-power conditions. The most cited is the
\textbf{LaSalle Unit~2 event of 1988}, in which a loss of feedwater heating and a
trip of both recirculation pumps left the reactor at high power and low (natural
circulation) flow---exactly the regime of weakest stability margin. The reactor
entered a growing power oscillation (an increasing decay ratio) before operators
manually scrammed it. No damage resulted, but the event prompted industry-wide
changes: stability exclusion regions on the power--flow map, oscillation detection
systems, and revised procedures. LaSalle is the clearest modern illustration of a
\emph{linearly} unstable operating point in a commercial reactor---and, unlike
Chernobyl, of a system caught and shut down before harm.

\section{Three Mile Island (1979): the thermal aftermath}

The Three Mile Island Unit~2 accident was not a reactivity-stability event---the
reactor scrammed correctly and promptly. It belongs in this history for the
\emph{other} half of thermal safety: decay heat. A stuck-open relief valve and
misread instruments led operators to throttle emergency cooling, and decay heat
(Chapter~\ref{ch:decay}) slowly uncovered and melted much of the core over hours.
TMI is the canonical reminder that stopping fission does not stop the heat, and
that human factors and instrumentation are part of the safety system. Its lessons
about operator training, control-room design, and instrumentation directly
informed the international response to Chernobyl seven years later.

\section{What the history teaches}

Read together, these events sketch a coherent set of principles:
\begin{itemize}[nosep]
\item A single large, fast reactivity insertion can take a reactor prompt-critical
faster than any intervention (SL-1).
\item In a reactor with prompt \emph{negative} feedback, excursions are
self-limiting (BORAX, SPERT)---this is the design goal.
\item Phase-lagged feedback can make a nominally negative-feedback reactor
oscillate (LaSalle); margins and monitoring are required.
\item Decay heat makes loss of cooling dangerous even after shutdown (TMI,
Fermi-1).
\item Shutdown systems must be fast, reliable, and unambiguous (NRX).
\end{itemize}
The RBMK violated or weakened several of these at once: it had a positive feedback
mechanism (void), it was operated into its least stable regime, its shutdown
system had a fatal flaw, and its operators acted on an incomplete understanding of
the machine's physics. The next chapters examine how.


\section*{Exercises}
\addcontentsline{toc}{section}{Exercises}
\begin{enumerate}[leftmargin=*,itemsep=2pt]
\item What did the BORAX and SPERT tests demonstrate about under-moderated water reactors?
\item Summarise the SL-1 accident and the two general lessons it shares with Chernobyl.
\item Explain why the LaSalle-2 event is an example of a linearly unstable operating point.
\item Why is Three Mile Island classified here as a decay-heat rather than a stability accident?
\item Extract five general safety principles from the historical record.
\end{enumerate}
